% Options for packages loaded elsewhere
\PassOptionsToPackage{unicode}{hyperref}
\PassOptionsToPackage{hyphens}{url}
\PassOptionsToPackage{dvipsnames,svgnames,x11names}{xcolor}
%
\documentclass[
  10pt,
]{article}
\usepackage{amsmath,amssymb}
\usepackage{iftex}
\ifPDFTeX
  \usepackage[T1]{fontenc}
  \usepackage[utf8]{inputenc}
  \usepackage{textcomp} % provide euro and other symbols
\else % if luatex or xetex
  \usepackage{unicode-math} % this also loads fontspec
  \defaultfontfeatures{Scale=MatchLowercase}
  \defaultfontfeatures[\rmfamily]{Ligatures=TeX,Scale=1}
\fi
\usepackage{lmodern}
\ifPDFTeX\else
  % xetex/luatex font selection
  \setmainfont[]{DejaVu Sans}
  \setmonofont[]{DejaVu Sans Mono}
  \setmathfont[]{DejaVu Sans}
\fi
% Use upquote if available, for straight quotes in verbatim environments
\IfFileExists{upquote.sty}{\usepackage{upquote}}{}
\IfFileExists{microtype.sty}{% use microtype if available
  \usepackage[]{microtype}
  \UseMicrotypeSet[protrusion]{basicmath} % disable protrusion for tt fonts
}{}
\makeatletter
\@ifundefined{KOMAClassName}{% if non-KOMA class
  \IfFileExists{parskip.sty}{%
    \usepackage{parskip}
  }{% else
    \setlength{\parindent}{0pt}
    \setlength{\parskip}{6pt plus 2pt minus 1pt}}
}{% if KOMA class
  \KOMAoptions{parskip=half}}
\makeatother
\usepackage{xcolor}
\usepackage[margin=2.2cm]{geometry}
\usepackage{color}
\usepackage{fancyvrb}
\newcommand{\VerbBar}{|}
\newcommand{\VERB}{\Verb[commandchars=\\\{\}]}
\DefineVerbatimEnvironment{Highlighting}{Verbatim}{commandchars=\\\{\}}
% Add ',fontsize=\small' for more characters per line
\newenvironment{Shaded}{}{}
\newcommand{\AlertTok}[1]{\textcolor[rgb]{1.00,0.00,0.00}{\textbf{#1}}}
\newcommand{\AnnotationTok}[1]{\textcolor[rgb]{0.38,0.63,0.69}{\textbf{\textit{#1}}}}
\newcommand{\AttributeTok}[1]{\textcolor[rgb]{0.49,0.56,0.16}{#1}}
\newcommand{\BaseNTok}[1]{\textcolor[rgb]{0.25,0.63,0.44}{#1}}
\newcommand{\BuiltInTok}[1]{\textcolor[rgb]{0.00,0.50,0.00}{#1}}
\newcommand{\CharTok}[1]{\textcolor[rgb]{0.25,0.44,0.63}{#1}}
\newcommand{\CommentTok}[1]{\textcolor[rgb]{0.38,0.63,0.69}{\textit{#1}}}
\newcommand{\CommentVarTok}[1]{\textcolor[rgb]{0.38,0.63,0.69}{\textbf{\textit{#1}}}}
\newcommand{\ConstantTok}[1]{\textcolor[rgb]{0.53,0.00,0.00}{#1}}
\newcommand{\ControlFlowTok}[1]{\textcolor[rgb]{0.00,0.44,0.13}{\textbf{#1}}}
\newcommand{\DataTypeTok}[1]{\textcolor[rgb]{0.56,0.13,0.00}{#1}}
\newcommand{\DecValTok}[1]{\textcolor[rgb]{0.25,0.63,0.44}{#1}}
\newcommand{\DocumentationTok}[1]{\textcolor[rgb]{0.73,0.13,0.13}{\textit{#1}}}
\newcommand{\ErrorTok}[1]{\textcolor[rgb]{1.00,0.00,0.00}{\textbf{#1}}}
\newcommand{\ExtensionTok}[1]{#1}
\newcommand{\FloatTok}[1]{\textcolor[rgb]{0.25,0.63,0.44}{#1}}
\newcommand{\FunctionTok}[1]{\textcolor[rgb]{0.02,0.16,0.49}{#1}}
\newcommand{\ImportTok}[1]{\textcolor[rgb]{0.00,0.50,0.00}{\textbf{#1}}}
\newcommand{\InformationTok}[1]{\textcolor[rgb]{0.38,0.63,0.69}{\textbf{\textit{#1}}}}
\newcommand{\KeywordTok}[1]{\textcolor[rgb]{0.00,0.44,0.13}{\textbf{#1}}}
\newcommand{\NormalTok}[1]{#1}
\newcommand{\OperatorTok}[1]{\textcolor[rgb]{0.40,0.40,0.40}{#1}}
\newcommand{\OtherTok}[1]{\textcolor[rgb]{0.00,0.44,0.13}{#1}}
\newcommand{\PreprocessorTok}[1]{\textcolor[rgb]{0.74,0.48,0.00}{#1}}
\newcommand{\RegionMarkerTok}[1]{#1}
\newcommand{\SpecialCharTok}[1]{\textcolor[rgb]{0.25,0.44,0.63}{#1}}
\newcommand{\SpecialStringTok}[1]{\textcolor[rgb]{0.73,0.40,0.53}{#1}}
\newcommand{\StringTok}[1]{\textcolor[rgb]{0.25,0.44,0.63}{#1}}
\newcommand{\VariableTok}[1]{\textcolor[rgb]{0.10,0.09,0.49}{#1}}
\newcommand{\VerbatimStringTok}[1]{\textcolor[rgb]{0.25,0.44,0.63}{#1}}
\newcommand{\WarningTok}[1]{\textcolor[rgb]{0.38,0.63,0.69}{\textbf{\textit{#1}}}}
\usepackage{longtable,booktabs,array}
\usepackage{calc} % for calculating minipage widths
% Correct order of tables after \paragraph or \subparagraph
\usepackage{etoolbox}
\makeatletter
\patchcmd\longtable{\par}{\if@noskipsec\mbox{}\fi\par}{}{}
\makeatother
% Allow footnotes in longtable head/foot
\IfFileExists{footnotehyper.sty}{\usepackage{footnotehyper}}{\usepackage{footnote}}
\makesavenoteenv{longtable}
\setlength{\emergencystretch}{3em} % prevent overfull lines
\providecommand{\tightlist}{%
  \setlength{\itemsep}{0pt}\setlength{\parskip}{0pt}}
\setcounter{secnumdepth}{-\maxdimen} % remove section numbering

\usepackage{xcolor}
\definecolor{plblue}{HTML}{1F3864}
\usepackage{sectsty}
\allsectionsfont{\sffamily\color{plblue}}
\usepackage{fvextra}
\DefineVerbatimEnvironment{Highlighting}{Verbatim}{breaklines,commandchars=\\\{\},fontsize=\footnotesize}
\usepackage{etoolbox}
\AtBeginEnvironment{longtable}{\footnotesize}
\ifLuaTeX
  \usepackage{selnolig}  % disable illegal ligatures
\fi
\IfFileExists{bookmark.sty}{\usepackage{bookmark}}{\usepackage{hyperref}}
\IfFileExists{xurl.sty}{\usepackage{xurl}}{} % add URL line breaks if available
\urlstyle{same}
\hypersetup{
  colorlinks=true,
  linkcolor={Maroon},
  filecolor={Maroon},
  citecolor={Blue},
  urlcolor={Blue},
  pdfcreator={LaTeX via pandoc}}

\author{}
\date{}

\begin{document}

\hypertarget{document-2-benchmarks}{%
\section{Document 2: Benchmarks}\label{document-2-benchmarks}}

\textbf{Kit:} qae\_european 1.0.0 · \textbf{Scope:} the classical
methods amplitude estimation is compared with: which ones, why, how they
are implemented, with which parameters, and which alternatives were
considered and not used.

Companion documents: \href{circuit.md}{Document 1: Circuit} ·
\href{tests.md}{Document 3: Tests} · \href{../README.md}{User guide}.

\emph{This file is generated by \texttt{docs/build\_benchmarks\_doc.py}
from \texttt{docs/benchmarks\_head.md} and the source files, so the code
in section 6 is always the code that ran. Edit the head file or the
code, not this file.}

\begin{center}\rule{0.5\linewidth}{0.5pt}\end{center}

\hypertarget{1-what-a-benchmark-must-answer}{%
\subsection{1. What a benchmark must
answer}\label{1-what-a-benchmark-must-answer}}

The method under test reads a price from an amplitude that it estimates
with a number of \textbf{calls} to the circuit A. A benchmark is useful
if it answers one of three questions:

\begin{enumerate}
\def\labelenumi{\arabic{enumi}.}
\tightlist
\item
  \textbf{What is the right answer?} A reference price with no sampling
  error.
\item
  \textbf{Is the quantum estimator better than the classical sampler it
  is meant to replace,} measured on the same target and with the same
  number of calls?
\item
  \textbf{How does it compare with the best classical sampling method}
  for this problem, not just the plain one?
\end{enumerate}

\hypertarget{2-how-performance-is-measured}{%
\subsection{2. How performance is
measured}\label{2-how-performance-is-measured}}

The full definition is in Document 3, section 1.1. In short:

\begin{itemize}
\tightlist
\item
  \textbf{Performance is measured in calls}, not in time:

  \begin{itemize}
  \tightlist
  \item
    for amplitude estimation, one application of A;
  \item
    for Monte Carlo, one simulated path;
  \item
    for quasi-Monte Carlo, one point of the sequence.
  \end{itemize}
\item
  \textbf{Error is the root-mean-square error (RMSE)} over independent
  repetitions, against a stated target. The median absolute error is
  also reported.
\item
  \textbf{Two methods are compared by their relative efficiency} E(ε) =
  N\_comparator(ε) / N\_QAE(ε), the ratio of calls each needs to reach
  the same RMSE ε.
\end{itemize}

Wall-clock time is deliberately out of scope. A classical simulation of
a quantum circuit is far slower than Monte Carlo on the same machine,
and a timing of the simulator says nothing about hardware.

\hypertarget{3-the-benchmarks-chosen}{%
\subsection{3. The benchmarks chosen}\label{3-the-benchmarks-chosen}}

\begin{longtable}[]{@{}
  >{\raggedright\arraybackslash}p{(\columnwidth - 6\tabcolsep) * \real{0.2514}}
  >{\raggedright\arraybackslash}p{(\columnwidth - 6\tabcolsep) * \real{0.1725}}
  >{\raggedright\arraybackslash}p{(\columnwidth - 6\tabcolsep) * \real{0.3449}}
  >{\raggedright\arraybackslash}p{(\columnwidth - 6\tabcolsep) * \real{0.2112}}@{}}
\toprule\noalign{}
\begin{minipage}[b]{\linewidth}\raggedright
Benchmark
\end{minipage} & \begin{minipage}[b]{\linewidth}\raggedright
Target
\end{minipage} & \begin{minipage}[b]{\linewidth}\raggedright
Question it answers
\end{minipage} & \begin{minipage}[b]{\linewidth}\raggedright
Class
\end{minipage} \\
\midrule\noalign{}
\endhead
\bottomrule\noalign{}
\endlastfoot
Black-Scholes closed form & Continuous price & 1. The reference &
\texttt{BlackScholes} \\
Plain Monte Carlo on the grid & Grid price & 2. Like-for-like comparator
& \texttt{MonteCarlo("grid",\ n)} \\
Plain Monte Carlo, continuous & Continuous price & 2. What a desk would
get by plain sampling, including all errors &
\texttt{MonteCarlo("continuous")} \\
Quasi-Monte Carlo, scrambled Sobol & Continuous price & 3. Best
classical sampling method & \texttt{QuasiMonteCarlo} \\
\end{longtable}

All of them share one interface, \texttt{price(option,\ budget,\ rng)},
which returns an \texttt{Estimate(value,\ calls,\ stderr)}.

\hypertarget{31-black-scholes-closed-form-the-reference}{%
\subsubsection{3.1 Black-Scholes closed form: the
reference}\label{31-black-scholes-closed-form-the-reference}}

C = S₀ N(d₁) − K e\^{}(−rT) N(d₂), with d₁ = {[}ln(S₀/K) + (r + σ²/2)
T{]} / (σ√T) and d₂ = d₁ − σ√T.

\begin{itemize}
\tightlist
\item
  \textbf{Why.} The model is Black-Scholes, so the closed form is exact:
  every error against it is a true error of the method. It uses no
  samples (calls = 0).
\item
  \textbf{Accuracy.} \texttt{scipy.stats.norm.cdf} is accurate to
  machine precision over the campaign\textquotesingle s ranges. The
  smallest in-scope price is 0.01, well above any rounding issue.
\item
  \textbf{Limit.} It prices the continuous distribution, not the grid
  the circuit sees. Errors against it mix the discretisation bias and
  the statistical error. That is why the like-for-like comparator below
  targets the grid.
\end{itemize}

\hypertarget{32-plain-monte-carlo-on-the-grid-the-like-for-like-comparator}{%
\subsubsection{3.2 Plain Monte Carlo on the grid: the like-for-like
comparator}\label{32-plain-monte-carlo-on-the-grid-the-like-for-like-comparator}}

\begin{itemize}
\tightlist
\item
  \textbf{How.} It draws grid points i with the probabilities p\_i of
  the circuit\textquotesingle s \texttt{LognormalGrid}, and averages the
  discounted payoffs e\^{}(−rT) max(x\_i − K, 0). Drawing N points is
  implemented as one multinomial draw of the counts, which is equivalent
  and faster. The standard error is the sample standard deviation
  divided by √N.
\item
  \textbf{Why.} It estimates exactly the quantity the circuit encodes,
  the grid price, with the classical procedure the quantum algorithm
  replaces: sampling the distribution and averaging the payoff. It is a
  \textbf{stronger} comparator than measuring the circuit A directly
  without Grover steps. Direct measurement gives Bernoulli draws with
  probability a, scaled by g\_max: their variance g\_max² a(1 − a) is
  never below the variance of the payoff on the grid, and is 4.7 times
  larger for the reference option. Beating grid Monte Carlo is therefore
  the harder test. It isolates the claim that matters: amplitude
  estimation against the classical procedure it replaces. There is no
  discretisation bias and no difference in target, only statistical
  error against statistical error, at the same number of calls.
\item
  \textbf{Used for:} H1 and H2 (Document 3).
\end{itemize}

\hypertarget{33-plain-monte-carlo-continuous}{%
\subsubsection{3.3 Plain Monte Carlo,
continuous}\label{33-plain-monte-carlo-continuous}}

\begin{itemize}
\tightlist
\item
  \textbf{How.} S\_T = S₀ exp((r − σ²/2) T + σ√T Z), with Z standard
  normal, then the discounted average payoff. Unbiased for the
  Black-Scholes price.
\item
  \textbf{Why.} It is what a desk would actually run without a quantum
  computer. It shows the total error a user would see, against the
  continuous price, with no discretisation.
\end{itemize}

\hypertarget{34-quasi-monte-carlo-scrambled-sobol-the-best-classical-sampler}{%
\subsubsection{3.4 Quasi-Monte Carlo, scrambled Sobol: the best
classical
sampler}\label{34-quasi-monte-carlo-scrambled-sobol-the-best-classical-sampler}}

\begin{itemize}
\item
  \textbf{How.}

  \begin{itemize}
  \tightlist
  \item
    Take a one-dimensional Sobol sequence with Owen scrambling
    (\texttt{scipy.stats.qmc.Sobol}, \texttt{scramble=True}).
  \item
    Map each point u\_j to Z\_j = Φ⁻¹(u\_j), then to S\_T as in plain
    Monte Carlo, and average the discounted payoffs.
  \item
    The number of points is the largest power of two not above the
    budget, because Sobol\textquotesingle s balance properties hold for
    powers of two.
  \item
    Each repetition uses a fresh scramble, so repetitions are
    independent randomisations and their spread measures the error. The
    method reports no standard error of its own.
  \end{itemize}
\item
  \textbf{Why.} The payoff is a smooth function of one variable, apart
  from the kink at the strike. In one dimension, scrambled Sobol
  converges at close to 1/N, against 1/√N for plain Monte Carlo. That is
  the same rate amplitude estimation promises. Leaving it out would
  compare the quantum method only with the weakest classical sampler.
  Document 3 states in advance (H4) that amplitude estimation is
  \textbf{not} expected to beat it.
\item
  \textbf{Cost of the power-of-two rule.} At each budget, quasi-Monte
  Carlo uses fewer points than the others:

  \begin{longtable}[]{@{}
    >{\raggedright\arraybackslash}p{(\columnwidth - 10\tabcolsep) * \real{0.2389}}
    >{\raggedright\arraybackslash}p{(\columnwidth - 10\tabcolsep) * \real{0.1428}}
    >{\raggedright\arraybackslash}p{(\columnwidth - 10\tabcolsep) * \real{0.1428}}
    >{\raggedright\arraybackslash}p{(\columnwidth - 10\tabcolsep) * \real{0.1428}}
    >{\raggedright\arraybackslash}p{(\columnwidth - 10\tabcolsep) * \real{0.1564}}
    >{\raggedright\arraybackslash}p{(\columnwidth - 10\tabcolsep) * \real{0.1564}}@{}}
  \toprule\noalign{}
  \begin{minipage}[b]{\linewidth}\raggedright
  Budget (calls)
  \end{minipage} & \begin{minipage}[b]{\linewidth}\raggedright
  1,800
  \end{minipage} & \begin{minipage}[b]{\linewidth}\raggedright
  3,500
  \end{minipage} & \begin{minipage}[b]{\linewidth}\raggedright
  6,800
  \end{minipage} & \begin{minipage}[b]{\linewidth}\raggedright
  13,300
  \end{minipage} & \begin{minipage}[b]{\linewidth}\raggedright
  26,200
  \end{minipage} \\
  \midrule\noalign{}
  \endhead
  \bottomrule\noalign{}
  \endlastfoot
  Sobol points & 1,024 & 2,048 & 4,096 & 8,192 & 16,384 \\
  \end{longtable}

  That is between 57\% and 63\% of the budget. This handicaps
  quasi-Monte Carlo, never amplitude estimation. The number of points
  actually used is recorded (\texttt{calls\_qmc}), and relative
  efficiencies are computed with it.
\end{itemize}

\hypertarget{4-variance-reduction-none-in-the-comparators}{%
\subsection{4. Variance reduction: none in the
comparators}\label{4-variance-reduction-none-in-the-comparators}}

Plain Monte Carlo is run \textbf{without} variance reduction (no
antithetic variates, control variates, importance sampling or
stratification). Three reasons:

\begin{enumerate}
\def\labelenumi{\arabic{enumi}.}
\tightlist
\item
  \textbf{It is what amplitude estimation replaces.} The quadratic
  speed-up is stated against plain sampling. Variance reduction lowers
  the constant in front of 1/√N but not the rate.
\item
  \textbf{It applies to both sides.} The same ideas exist for amplitude
  estimation (for example, encoding a control variate in the payoff, or
  a better-adapted grid). Adding them to one side only would mix two
  effects.
\item
  \textbf{The best classical method is covered separately.} Scrambled
  Sobol (section 3.4) is itself a variance-reduction method, and the
  strongest one for this one-dimensional problem. It answers "what if
  the classical side tries hard".
\end{enumerate}

A reader who wants Monte Carlo with variance reduction can subclass
\texttt{Benchmark} (section 6) and run it through the same campaign
code.

\hypertarget{5-benchmarks-considered-and-not-used}{%
\subsection{5. Benchmarks considered and not
used}\label{5-benchmarks-considered-and-not-used}}

\begin{longtable}[]{@{}
  >{\raggedright\arraybackslash}p{(\columnwidth - 2\tabcolsep) * \real{0.3020}}
  >{\raggedright\arraybackslash}p{(\columnwidth - 2\tabcolsep) * \real{0.6780}}@{}}
\toprule\noalign{}
\begin{minipage}[b]{\linewidth}\raggedright
Candidate
\end{minipage} & \begin{minipage}[b]{\linewidth}\raggedright
Why not used
\end{minipage} \\
\midrule\noalign{}
\endhead
\bottomrule\noalign{}
\endlastfoot
\textbf{Deterministic quadrature} (e.g. Gauss-Hermite, or the exact sum
over the grid) & For a one-dimensional integral, quadrature beats every
sampling method: it is exact to machine precision in a few dozen points.
It is used, implicitly, as the reference: the grid price is that exact
sum, and Black-Scholes is the exact integral. As a \emph{comparator} it
would only state that sampling is the wrong tool in one dimension. That
is true and is said here, but amplitude estimation is meant for problems
where quadrature fails (many dimensions, path dependence), and this kit
only tests the method on the simplest case. \\
\textbf{Monte Carlo with variance reduction} & See section 4: it changes
the constant, not the rate, and would have to be matched on the quantum
side. Left to users through the \texttt{Benchmark} interface. \\
\textbf{Multilevel Monte Carlo} & It reduces the cost of time
discretisation in path simulation. Here S\_T is sampled exactly in one
step, so there is nothing for it to reduce. \\
\textbf{Other amplitude-estimation variants} (canonical QAE with phase
estimation, iterative QAE, power-law MLAE) & These are alternatives to
the method under test, not classical benchmarks. Canonical QAE needs
controlled Grover operators and deeper circuits. Iterative QAE is a
natural next comparison, out of scope for the simplest kit. \\
\textbf{An external quantum implementation} (e.g.
Qiskit\textquotesingle s finance application module) & It would
cross-check the circuit against an independent implementation of the
same construction. Not used, to keep the dependencies to Qiskit, NumPy
and SciPy. The circuit is instead checked against the exact grid price
(H0-a) and by full statevector simulation (H0-b). \\
\textbf{Wall-clock timing} & Out of scope (section 2). \\
\end{longtable}

\hypertarget{6-the-code}{%
\subsection{6. The code}\label{6-the-code}}

The benchmarks live in \texttt{src/qae\_european/benchmarks/}. Below are
the full source files, comments included, exactly as run.

\hypertarget{61-basepy-shared-interface-and-result-type}{%
\subsubsection{\texorpdfstring{6.1 \texttt{base.py}: Shared interface
and result
type}{6.1 base.py: Shared interface and result type}}\label{61-basepy-shared-interface-and-result-type}}

\begin{Shaded}
\begin{Highlighting}[]
\CommentTok{"""Shared engine of the benchmarks: the result type and the common interface."""}
\ImportTok{from}\NormalTok{ dataclasses }\ImportTok{import}\NormalTok{ dataclass}
\ImportTok{import}\NormalTok{ numpy }\ImportTok{as}\NormalTok{ np}


\AttributeTok{@dataclass}\NormalTok{(frozen}\OperatorTok{=}\VariableTok{True}\NormalTok{)}
\KeywordTok{class}\NormalTok{ Estimate:}
    \CommentTok{"""Result of one pricing run.}

\CommentTok{    value : float         price estimate}
\CommentTok{    calls : int           number of payoff evaluations (paths or points) actually used; 0 for a closed form}
\CommentTok{    stderr : float        standard error reported by the method itself, if it has one (NaN otherwise)}
\CommentTok{    """}
\NormalTok{    value: }\BuiltInTok{float}
\NormalTok{    calls: }\BuiltInTok{int}
\NormalTok{    stderr: }\BuiltInTok{float} \OperatorTok{=} \BuiltInTok{float}\NormalTok{(}\StringTok{"nan"}\NormalTok{)}


\KeywordTok{class}\NormalTok{ Benchmark:}
    \CommentTok{"""Interface: price(option, budget, rng) returns an Estimate using at most \textasciigrave{}budget\textasciigrave{} payoff evaluations."""}
\NormalTok{    name }\OperatorTok{=} \StringTok{"benchmark"}

    \KeywordTok{def}\NormalTok{ price(}\VariableTok{self}\NormalTok{, option, budget: }\BuiltInTok{int}\NormalTok{, rng: np.random.Generator) }\OperatorTok{{-}\textgreater{}}\NormalTok{ Estimate:  }\CommentTok{\# pragma: no cover}
        \ControlFlowTok{raise} \PreprocessorTok{NotImplementedError}

    \AttributeTok{@staticmethod}
    \KeywordTok{def}\NormalTok{ \_discounted\_payoffs(option, terminal\_prices: np.ndarray) }\OperatorTok{{-}\textgreater{}}\NormalTok{ np.ndarray:}
        \CommentTok{"""Shared step: discounted call payoffs for an array of terminal prices."""}
        \ControlFlowTok{return}\NormalTok{ option.discount }\OperatorTok{*}\NormalTok{ option.payoff(terminal\_prices)}
\end{Highlighting}
\end{Shaded}

\hypertarget{62-black_scholespy-black-scholes-closed-form}{%
\subsubsection{\texorpdfstring{6.2 \texttt{black\_scholes.py}:
Black-Scholes closed
form}{6.2 black\_scholes.py: Black-Scholes closed form}}\label{62-black_scholespy-black-scholes-closed-form}}

\begin{Shaded}
\begin{Highlighting}[]
\CommentTok{"""Black{-}Scholes closed form for a European call: the reference price."""}
\ImportTok{import}\NormalTok{ math}
\ImportTok{from}\NormalTok{ scipy.stats }\ImportTok{import}\NormalTok{ norm}

\ImportTok{from}\NormalTok{ .base }\ImportTok{import}\NormalTok{ Benchmark, Estimate}


\KeywordTok{class}\NormalTok{ BlackScholes(Benchmark):}
    \CommentTok{"""C = S0 N(d1) {-} K exp({-}rT) N(d2),}
\CommentTok{    d1 = [ln(S0/K) + (r + sigma\^{}2/2) T] / (sigma sqrt T),  d2 = d1 {-} sigma sqrt T.}

\CommentTok{    Exact under the model; uses no samples (calls = 0). Implemented with scipy.stats.norm.cdf,}
\CommentTok{    which is accurate to machine precision across the parameter ranges of the campaign.}
\CommentTok{    """}
\NormalTok{    name }\OperatorTok{=} \StringTok{"black\_scholes"}

    \AttributeTok{@staticmethod}
    \KeywordTok{def}\NormalTok{ call(option) }\OperatorTok{{-}\textgreater{}} \BuiltInTok{float}\NormalTok{:}
\NormalTok{        sq }\OperatorTok{=}\NormalTok{ option.log\_std}
\NormalTok{        d1 }\OperatorTok{=}\NormalTok{ (math.log(option.spot }\OperatorTok{/}\NormalTok{ option.strike) }\OperatorTok{+}\NormalTok{ (option.rate }\OperatorTok{+} \FloatTok{0.5} \OperatorTok{*}\NormalTok{ option.vol }\OperatorTok{**} \DecValTok{2}\NormalTok{) }\OperatorTok{*}\NormalTok{ option.maturity) }\OperatorTok{/}\NormalTok{ sq}
\NormalTok{        d2 }\OperatorTok{=}\NormalTok{ d1 }\OperatorTok{{-}}\NormalTok{ sq}
        \ControlFlowTok{return} \BuiltInTok{float}\NormalTok{(option.spot }\OperatorTok{*}\NormalTok{ norm.cdf(d1) }\OperatorTok{{-}}\NormalTok{ option.strike }\OperatorTok{*}\NormalTok{ option.discount }\OperatorTok{*}\NormalTok{ norm.cdf(d2))}

    \KeywordTok{def}\NormalTok{ price(}\VariableTok{self}\NormalTok{, option, budget: }\BuiltInTok{int} \OperatorTok{=} \DecValTok{0}\NormalTok{, rng}\OperatorTok{=}\VariableTok{None}\NormalTok{) }\OperatorTok{{-}\textgreater{}}\NormalTok{ Estimate:}
        \ControlFlowTok{return}\NormalTok{ Estimate(}\VariableTok{self}\NormalTok{.call(option), }\DecValTok{0}\NormalTok{, }\FloatTok{0.0}\NormalTok{)}
\end{Highlighting}
\end{Shaded}

\hypertarget{63-monte_carlopy-plain-monte-carlo-continuous-and-on-the-grid}{%
\subsubsection{\texorpdfstring{6.3 \texttt{monte\_carlo.py}: Plain Monte
Carlo, continuous and on the
grid}{6.3 monte\_carlo.py: Plain Monte Carlo, continuous and on the grid}}\label{63-monte_carlopy-plain-monte-carlo-continuous-and-on-the-grid}}

\begin{Shaded}
\begin{Highlighting}[]
\CommentTok{"""Plain Monte Carlo, with two sampling settings.}

\CommentTok{sampling="continuous": draws S\_T exactly from the lognormal, S\_T = S0 exp((r {-} sigma\^{}2/2) T + sigma sqrt(T) Z),}
\CommentTok{                       Z standard normal. Unbiased for the Black{-}Scholes price.}
\CommentTok{sampling="grid":       draws grid points with the probabilities of the circuit\textquotesingle{}s LognormalGrid}
\CommentTok{                       (n price qubits). Unbiased for the grid price, i.e. the same target as the circuit:}
\CommentTok{                       it isolates statistical error from discretisation bias.}

\CommentTok{No variance reduction: plain sampling is exactly what amplitude estimation replaces.}
\CommentTok{Standard error = sample standard deviation / sqrt(budget).}
\CommentTok{"""}
\ImportTok{import}\NormalTok{ numpy }\ImportTok{as}\NormalTok{ np}

\ImportTok{from}\NormalTok{ .base }\ImportTok{import}\NormalTok{ Benchmark, Estimate}
\ImportTok{from}\NormalTok{ ..grid }\ImportTok{import}\NormalTok{ LognormalGrid}


\KeywordTok{class}\NormalTok{ MonteCarlo(Benchmark):}
    \KeywordTok{def} \FunctionTok{\_\_init\_\_}\NormalTok{(}\VariableTok{self}\NormalTok{, sampling: }\BuiltInTok{str} \OperatorTok{=} \StringTok{"continuous"}\NormalTok{, n\_qubits: }\BuiltInTok{int} \OperatorTok{|} \VariableTok{None} \OperatorTok{=} \VariableTok{None}\NormalTok{):}
        \ControlFlowTok{if}\NormalTok{ sampling }\KeywordTok{not} \KeywordTok{in}\NormalTok{ (}\StringTok{"continuous"}\NormalTok{, }\StringTok{"grid"}\NormalTok{):}
            \ControlFlowTok{raise} \PreprocessorTok{ValueError}\NormalTok{(}\StringTok{"sampling must be \textquotesingle{}continuous\textquotesingle{} or \textquotesingle{}grid\textquotesingle{}"}\NormalTok{)}
        \ControlFlowTok{if}\NormalTok{ sampling }\OperatorTok{==} \StringTok{"grid"} \KeywordTok{and}\NormalTok{ n\_qubits }\KeywordTok{is} \VariableTok{None}\NormalTok{:}
            \ControlFlowTok{raise} \PreprocessorTok{ValueError}\NormalTok{(}\StringTok{"grid sampling needs n\_qubits"}\NormalTok{)}
        \VariableTok{self}\NormalTok{.sampling, }\VariableTok{self}\NormalTok{.n\_qubits }\OperatorTok{=}\NormalTok{ sampling, n\_qubits}
        \VariableTok{self}\NormalTok{.name }\OperatorTok{=} \StringTok{"monte\_carlo"} \ControlFlowTok{if}\NormalTok{ sampling }\OperatorTok{==} \StringTok{"continuous"} \ControlFlowTok{else} \SpecialStringTok{f"monte\_carlo\_grid\_}\SpecialCharTok{\{}\NormalTok{n\_qubits}\SpecialCharTok{\}}\SpecialStringTok{"}

    \KeywordTok{def}\NormalTok{ price(}\VariableTok{self}\NormalTok{, option, budget: }\BuiltInTok{int}\NormalTok{, rng: np.random.Generator) }\OperatorTok{{-}\textgreater{}}\NormalTok{ Estimate:}
        \ControlFlowTok{if}\NormalTok{ budget }\OperatorTok{\textless{}} \DecValTok{2}\NormalTok{:}
            \ControlFlowTok{raise} \PreprocessorTok{ValueError}\NormalTok{(}\StringTok{"budget must be at least 2"}\NormalTok{)}
        \ControlFlowTok{if} \VariableTok{self}\NormalTok{.sampling }\OperatorTok{==} \StringTok{"continuous"}\NormalTok{:}
\NormalTok{            z }\OperatorTok{=}\NormalTok{ rng.standard\_normal(budget)}
\NormalTok{            st }\OperatorTok{=}\NormalTok{ np.exp(option.log\_mean }\OperatorTok{+}\NormalTok{ option.log\_std }\OperatorTok{*}\NormalTok{ z)}
\NormalTok{            pay }\OperatorTok{=} \VariableTok{self}\NormalTok{.\_discounted\_payoffs(option, st)}
            \ControlFlowTok{return}\NormalTok{ Estimate(}\BuiltInTok{float}\NormalTok{(pay.mean()), budget, }\BuiltInTok{float}\NormalTok{(pay.std(ddof}\OperatorTok{=}\DecValTok{1}\NormalTok{) }\OperatorTok{/}\NormalTok{ np.sqrt(budget)))}
\NormalTok{        g }\OperatorTok{=}\NormalTok{ LognormalGrid(option, }\VariableTok{self}\NormalTok{.n\_qubits)}
\NormalTok{        counts }\OperatorTok{=}\NormalTok{ rng.multinomial(budget, g.probs)             }\CommentTok{\# equivalent to drawing \textasciigrave{}budget\textasciigrave{} grid points}
\NormalTok{        pay }\OperatorTok{=}\NormalTok{ option.discount }\OperatorTok{*}\NormalTok{ g.payoffs()}
\NormalTok{        mean }\OperatorTok{=} \BuiltInTok{float}\NormalTok{(np.dot(counts, pay) }\OperatorTok{/}\NormalTok{ budget)}
\NormalTok{        var }\OperatorTok{=} \BuiltInTok{float}\NormalTok{(np.dot(counts, (pay }\OperatorTok{{-}}\NormalTok{ mean) }\OperatorTok{**} \DecValTok{2}\NormalTok{) }\OperatorTok{/}\NormalTok{ (budget }\OperatorTok{{-}} \DecValTok{1}\NormalTok{))}
        \ControlFlowTok{return}\NormalTok{ Estimate(mean, budget, }\BuiltInTok{float}\NormalTok{(np.sqrt(var }\OperatorTok{/}\NormalTok{ budget)))}
\end{Highlighting}
\end{Shaded}

\hypertarget{64-quasi_monte_carlopy-quasi-monte-carlo-scrambled-sobol}{%
\subsubsection{\texorpdfstring{6.4 \texttt{quasi\_monte\_carlo.py}:
Quasi-Monte Carlo, scrambled
Sobol}{6.4 quasi\_monte\_carlo.py: Quasi-Monte Carlo, scrambled Sobol}}\label{64-quasi_monte_carlopy-quasi-monte-carlo-scrambled-sobol}}

\begin{Shaded}
\begin{Highlighting}[]
\CommentTok{"""Quasi{-}Monte Carlo with a scrambled Sobol sequence (Owen scrambling, scipy.stats.qmc.Sobol).}

\CommentTok{One{-}dimensional problem: u\_j in (0,1) from the Sobol sequence, Z\_j = Phi\^{}\{{-}1\}(u\_j), then S\_T as in plain}
\CommentTok{Monte Carlo. The number of points is the largest power of two not above the budget (Sobol\textquotesingle{}s balance}
\CommentTok{properties hold for powers of two), so \textasciigrave{}calls\textasciigrave{} can be below the budget; the campaign records it.}

\CommentTok{Each call uses a fresh scramble seeded from \textasciigrave{}rng\textasciigrave{}, so repeated runs are independent randomisations:}
\CommentTok{the campaign measures the error across repetitions. The method itself reports no standard error.}
\CommentTok{"""}
\ImportTok{import}\NormalTok{ numpy }\ImportTok{as}\NormalTok{ np}
\ImportTok{from}\NormalTok{ scipy.stats }\ImportTok{import}\NormalTok{ norm, qmc}

\ImportTok{from}\NormalTok{ .base }\ImportTok{import}\NormalTok{ Benchmark, Estimate}


\KeywordTok{class}\NormalTok{ QuasiMonteCarlo(Benchmark):}
\NormalTok{    name }\OperatorTok{=} \StringTok{"quasi\_monte\_carlo"}

    \KeywordTok{def}\NormalTok{ price(}\VariableTok{self}\NormalTok{, option, budget: }\BuiltInTok{int}\NormalTok{, rng: np.random.Generator) }\OperatorTok{{-}\textgreater{}}\NormalTok{ Estimate:}
        \ControlFlowTok{if}\NormalTok{ budget }\OperatorTok{\textless{}} \DecValTok{2}\NormalTok{:}
            \ControlFlowTok{raise} \PreprocessorTok{ValueError}\NormalTok{(}\StringTok{"budget must be at least 2"}\NormalTok{)}
\NormalTok{        m }\OperatorTok{=} \BuiltInTok{int}\NormalTok{(np.floor(np.log2(budget)))}
\NormalTok{        sobol }\OperatorTok{=}\NormalTok{ qmc.Sobol(d}\OperatorTok{=}\DecValTok{1}\NormalTok{, scramble}\OperatorTok{=}\VariableTok{True}\NormalTok{, seed}\OperatorTok{=}\NormalTok{rng)}
\NormalTok{        u }\OperatorTok{=}\NormalTok{ sobol.random\_base2(m).ravel()}
\NormalTok{        st }\OperatorTok{=}\NormalTok{ np.exp(option.log\_mean }\OperatorTok{+}\NormalTok{ option.log\_std }\OperatorTok{*}\NormalTok{ norm.ppf(u))}
\NormalTok{        pay }\OperatorTok{=} \VariableTok{self}\NormalTok{.\_discounted\_payoffs(option, st)}
        \ControlFlowTok{return}\NormalTok{ Estimate(}\BuiltInTok{float}\NormalTok{(pay.mean()), }\DecValTok{2} \OperatorTok{**}\NormalTok{ m)}
\end{Highlighting}
\end{Shaded}

\hypertarget{7-parameters-used-in-the-campaign}{%
\subsection{7. Parameters used in the
campaign}\label{7-parameters-used-in-the-campaign}}

\begin{longtable}[]{@{}
  >{\raggedright\arraybackslash}p{(\columnwidth - 4\tabcolsep) * \real{0.2685}}
  >{\raggedright\arraybackslash}p{(\columnwidth - 4\tabcolsep) * \real{0.4840}}
  >{\raggedright\arraybackslash}p{(\columnwidth - 4\tabcolsep) * \real{0.2276}}@{}}
\toprule\noalign{}
\begin{minipage}[b]{\linewidth}\raggedright
Parameter
\end{minipage} & \begin{minipage}[b]{\linewidth}\raggedright
Value
\end{minipage} & \begin{minipage}[b]{\linewidth}\raggedright
Where
\end{minipage} \\
\midrule\noalign{}
\endhead
\bottomrule\noalign{}
\endlastfoot
Budgets (calls) & 1,800 · 3,500 · 6,800 · 13,300 · 26,200 (equal to the
amplitude-estimation ladders of depth 4 to 8) & \texttt{campaign.py} \\
Sobol points & largest power of two not above the budget &
\texttt{quasi\_monte\_carlo.py} \\
Repetitions per configuration & 50, independent &
\texttt{campaign.DRAWS} \\
Seeds & master seed 20261103; one stream per (scenario, qubits, depth,
method) through \texttt{numpy.random.SeedSequence} &
\texttt{campaign.rng\_for} \\
Variance reduction & none in Monte Carlo; Owen scrambling in quasi-Monte
Carlo & section 4 \\
Grid for Monte Carlo on the grid & the circuit\textquotesingle s own
\texttt{LognormalGrid(option,\ n)}, width 4 standard deviations &
\texttt{monte\_carlo.py}, \texttt{grid.py} \\
\end{longtable}

Continuous Monte Carlo and quasi-Monte Carlo do not depend on the number
of qubits. They are computed once per scenario and budget, and repeated
on each qubit row of the results.

\hypertarget{8-unit-tests-of-the-benchmarks}{%
\subsection{8. Unit tests of the
benchmarks}\label{8-unit-tests-of-the-benchmarks}}

\texttt{tests/test\_package.py} checks:

\begin{itemize}
\tightlist
\item
  the Black-Scholes price against the textbook value (S₀ = K = 100, r =
  5\%, σ = 20\%, T = 1: 10.4506) to 10⁻⁹, and the no-arbitrage bounds on
  a second case;
\item
  that both Monte Carlo settings are unbiased: with 200,000 paths,
  within 4 standard errors of their own target (Black-Scholes, grid
  price);
\item
  that quasi-Monte Carlo uses 4,096 points for a budget of 5,000 and
  lands within 0.02 of Black-Scholes;
\item
  that a campaign row is identical when recomputed from its seed.
\end{itemize}

Run them with \texttt{pytest\ tests}.

\end{document}
